\begin{abstract}
Comparing candidate classifiers on group fairness requires true labels, and under simple random sampling such comparisons need far more labels than comparisons of accuracy, because group rates have small denominators and the contrasts of interest are differences of group gaps. We study how to spend a fixed labelling budget when several trained models are compared on group-specific true positive rates, false positive rates and error. Two structural facts are used: differences between models arise only where their predictions differ, which is rare, and group membership is known before labelling. The method, PATTERN-DIF, stratifies the pool by group and by the number of models predicting the positive class, allocates labels by a weighted generalised Neyman rule with the models' probabilities as proxies, and estimates all contrasts with a stratified difference estimator that uses the proxy means on the whole pool, with linearisation intervals. Disagreement-based sampling and model-assisted estimation are not new; our contribution is their combination for ratio-type fairness contrasts among several models, with theory and evidence. Over eleven data settings with nine models, the median efficiency relative to simple random sampling is 5.8, 6.3 and 6.9 at 250, 500 and 1000 labels, against 3.7, 3.9 and 4.4 for the plain estimator on the same design and about 1.0 for stratification by group alone. Most of the gain comes from the disagreement strata. Coverage is close to nominal on average but poor for the worst contrast in one setting, and evidence rests on single splits.

\medskip\noindent\textbf{Keywords:} algorithmic fairness; model comparison; stratified sampling; model-assisted estimation; label-efficient evaluation; equalised odds
\end{abstract}
